\section{標準模型への展望}

\noindent\textbf{この章の中心命題}
\[
\boxed{
G_{\mathrm{SM}}
=
SU(3)_c\times SU(2)_L\times U(1)_Y,
\qquad
Q=T_3+Y
}
\]
標準模型は、強い相互作用を $SU(3)_c$、弱い相互作用と電磁相互作用を $SU(2)_L\times U(1)_Y$ のゲージ理論としてまとめ、ヒッグス場によって弱い相互作用の媒介粒子に質量を与える理論である。

ここまでで、QED が局所 $U(1)$ ゲージ対称性から構成されること、そして一般のゲージ理論では非可換な内部対称性も扱えることを見た。標準模型は、この考え方を電磁相互作用だけでなく、強い相互作用と弱い相互作用にも適用した理論である。

この章は展望であり、標準模型の全てを導出することは目的にしない。目的は、これまでに導入した道具がどのように標準模型の構造へ接続するかを明確にすることである。

\subsection{強い相互作用と色}
\begin{experimentalfact}{ハドロンの内部構造とジェット}
\factitem{操作}{高エネルギーの電子や陽子を陽子に衝突させ、散乱後に出てくる粒子の角度分布やエネルギー分布を測定する。}
\factitem{結果}{陽子や中性子は点状粒子のようには振る舞わず、内部により小さな構成要素があるような散乱分布を示す。また高エネルギー衝突では、狭い方向に多数のハドロンがまとまって飛ぶジェットが観測される。}
\factitem{示唆}{ハドロンの内部にはクォークとグルーオンの自由度があり、強い相互作用はそれらを結び付ける理論として記述する必要がある。}
\end{experimentalfact}

強い相互作用を記述する理論は量子色力学（QCD）である。QCD では、クォークは色と呼ばれる三種類の内部自由度をもつ。
\[
q(x)=
\begin{pmatrix}
q_r(x)\\
q_g(x)\\
q_b(x)
\end{pmatrix}.
\]
ここで $r,g,b$ は赤・緑・青という名前で呼ばれるが、これは実際の色ではなく、三成分の内部自由度を区別する記号である。

色の基底を時空点ごとに取り替えても物理が変わらないと要請すると、ゲージ群は $SU(3)_c$ になる。対応するゲージ場は八種類のグルーオン場 $G_\mu^a$ であり、共変微分は
\[
D_\mu q
=
\left(\partial_\mu+ig_sG_\mu^aT^a\right)q
\qquad(a=1,\dots,8)
\]
である。$T^a$ は $SU(3)$ の生成子、$g_s$ は強い相互作用の結合定数である。

QCD の基本的なラグランジアンは
\[
\mathcal{L}_{\mathrm{QCD}}
=
\bar{q}(i\gamma^\mu D_\mu-m_q)q
-\frac{1}{4}G_{\mu\nu}^aG^{a\mu\nu}
\]
である。$SU(3)$ は非可換なので、場の強さには
\[
G_{\mu\nu}^a
=
\partial_\mu G_\nu^a-\partial_\nu G_\mu^a
-g_s f^{abc}G_\mu^bG_\nu^c
\]
という自己相互作用項が含まれる。したがって、グルーオンはクォーク同士を結び付けるだけでなく、グルーオン同士でも相互作用する。

QCD の特徴は、短距離では相互作用が弱くなり、高エネルギー散乱ではクォークやグルーオンがほぼ自由粒子のように振る舞う一方、長距離では相互作用が強くなって孤立したクォークが観測されないことである。前者を漸近的自由性、後者を閉じ込めという。

\subsection{弱い相互作用と左巻き二重項}
\begin{experimentalfact}{弱い相互作用の媒介粒子}
\factitem{操作}{高エネルギー衝突実験で、荷電レプトンとニュートリノ、または荷電レプトン対に崩壊する重い粒子を探索する。}
\factitem{結果}{電荷をもつ $W^\pm$ ボソンと、中性の $Z$ ボソンが観測される。}
\factitem{示唆}{弱い相互作用は接触的な力ではなく、重いゲージ粒子の交換として記述される。}
\end{experimentalfact}

弱い相互作用の大きな特徴は、左巻き成分と右巻き成分を同じようには扱わないことである。標準模型では、左巻きのレプトンを
\[
L=
\begin{pmatrix}
\nu_e\\
e
\end{pmatrix}_L
\]
のような $SU(2)_L$ 二重項としてまとめる。一方、右巻き電子 $e_R$ は $SU(2)_L$ の一重項として扱う。この非対称な割り当てが、弱い相互作用で $P$ が破れることと結び付いている。

$SU(2)_L$ のゲージ場を $W_\mu^i$、$U(1)_Y$ のゲージ場を $B_\mu$ と書くと、電弱理論の共変微分は模式的に
\[
D_\mu
=
\partial_\mu
+ig W_\mu^i\frac{\tau^i}{2}
+ig'Y B_\mu
\]
である。ここで $\tau^i$ はパウリ行列、$Y$ はハイパーチャージ、$g,g'$ は結合定数である。電荷は
\[
Q=T_3+Y
\]
で与えられる。ここで $T_3$ は $SU(2)_L$ の第三成分であり、本ノートではこの式が成り立つように $Y$ の規格化を選んでいる。

\subsection{ヒッグス機構と質量}
ゲージ対称性をそのまま保つだけでは、$W^\pm$ や $Z$ に質量項
\[
\frac{1}{2}M^2 W_\mu W^\mu
\]
を手で足すことはできない。この項は局所ゲージ変換の下で不変ではないからである。しかし実験では、弱い相互作用の媒介粒子は大きな質量をもつ。したがって、ゲージ対称性を基礎に保ったまま、物理的には質量をもつ粒子が現れる仕組みが必要になる。

標準模型では、複素二重項のヒッグス場
\[
\Phi=
\begin{pmatrix}
\phi^+\\
\phi^0
\end{pmatrix}
\]
を導入し、ポテンシャルを
\[
V(\Phi)
=
-\mu^2\Phi^\dagger\Phi
+\lambda(\Phi^\dagger\Phi)^2
\qquad(\mu^2>0,\ \lambda>0)
\]
とする。このポテンシャルの最小値は $\Phi=0$ ではなく、
\[
\langle\Phi\rangle
=
\frac{1}{\sqrt{2}}
\begin{pmatrix}
0\\
v
\end{pmatrix}
\]
で与えられる。この真空期待値により、
\[
SU(2)_L\times U(1)_Y
\longrightarrow
U(1)_{\mathrm{em}}
\]
という形で対称性が自発的に破れる。

その結果、三つのゲージ場の組み合わせが質量をもつ $W^\pm,Z$ になり、一つの組み合わせが質量ゼロの光子 $A_\mu$ として残る。質量は
\[
M_W=\frac{1}{2}gv,
\qquad
M_Z=\frac{1}{2}\sqrt{g^2+g'^2}\,v
\]
となる。電磁相互作用が長距離力として残るのは、破れずに残った対称性が $U(1)_{\mathrm{em}}$ だからである。

\subsection{標準模型の全体像}
標準模型のゲージ群は
\[
G_{\mathrm{SM}}
=
SU(3)_c\times SU(2)_L\times U(1)_Y
\]
である。対応するゲージ場は、強い相互作用のグルーオン $G_\mu^a$、弱い相互作用の $W_\mu^i$、ハイパーチャージの $B_\mu$ である。物質場はクォークとレプトンであり、それぞれの左巻き成分・右巻き成分がこのゲージ群のもとで異なる変換性をもつ。

ラグランジアンを模式的に書けば、
\[
\mathcal{L}_{\mathrm{SM}}
=
\mathcal{L}_{\mathrm{gauge}}
+\mathcal{L}_{\mathrm{fermion}}
+\mathcal{L}_{\mathrm{Higgs}}
+\mathcal{L}_{\mathrm{Yukawa}}
\]
である。各項の役割は次の通りである。
\begin{itemize}
\item $\mathcal{L}_{\mathrm{gauge}}$ は、ゲージ場の運動を表す。
\item $\mathcal{L}_{\mathrm{fermion}}$ は、クォーク・レプトンの運動とゲージ相互作用を表す。
\item $\mathcal{L}_{\mathrm{Higgs}}$ は、ヒッグス場と対称性の自発的破れを表す。
\item $\mathcal{L}_{\mathrm{Yukawa}}$ は、フェルミオン質量の起源を表す。
\end{itemize}

このように標準模型は、三つの相互作用を単に並べた表ではない。基本原理は、局所ゲージ対称性と場の量子論である。QED で見た「局所対称性が相互作用の形を決める」という考え方が、非可換ゲージ理論、繰り込み、離散対称性の実験事実を経て、強い相互作用と弱い相互作用を含む統一的な枠組みに拡張されている。
